\hypertarget{phuxe1t-triux1ec3n-crawler}{%
\section{Phát triển Crawler}}

Phần này bao gồm việc xây dựng crawler Scrapy để thu thập bài báo từ các
trang tin tức Việt Nam sử dụng SitemapSpider.

\hypertarget{tux1ea1i-sao-duxf9ng-sitemapspider}{%
\subsection{Tại sao dùng
SitemapSpider?}}

Khác với \texttt{CrawlSpider} truyền thống dùng \texttt{DEPTH\_LIMIT},
\texttt{SitemapSpider} đọc trực tiếp \texttt{sitemap\_news.xml}, đảm
bảo: - \textbf{Không bỏ sót bài} --- phát hiện tất cả URL trong sitemap,
không chỉ các trang link - \textbf{Truy cập lịch sử} --- có thể crawl
bài cũ qua sitemap archives - \textbf{Phát hiện nhanh hơn} --- không cần
follow pagination links - \textbf{Tuân thủ robots.txt} --- built-in
\texttt{ROBOTSTXT\_OBEY\ =\ True}

\hypertarget{cux1ea5u-truxfac-dux1ef1-uxe1n}{%
\subsection{Cấu trúc dự án}}

\begin{verbatim}
crawler/
|--- __init__.py
|--- settings.py           # Cài đặt Scrapy
|--- pipelines.py          # Item pipelines (Kafka, PostgreSQL)
|--- items.py              # Data models
`--- spiders/
    |--- __init__.py
    `--- spider.py         # SitemapSpider implementation
\end{verbatim}

\hypertarget{cux1ea5u-huxecnh-configconfig_site.json}{%
\subsection{Cấu hình
(config/config\_site.json)}}

\emph{(Mã nguồn chi tiết đã được lược bỏ để báo cáo ngắn gọn. Vui lòng
xem trong source code dự án)}

Mỗi trang định nghĩa XPath selectors để trích xuất dữ liệu có cấu trúc.

\hypertarget{cuxe0i-ux111ux1eb7t-scrapy-crawlersettings.py}{%
\subsection{Cài đặt Scrapy
(crawler/settings.py)}}

\emph{(Mã nguồn chi tiết đã được lược bỏ để báo cáo ngắn gọn. Vui lòng
xem trong source code dự án)}

\hypertarget{items-crawleritems.py}{%
\subsection{Items (crawler/items.py)}}

\emph{(Mã nguồn chi tiết đã được lược bỏ để báo cáo ngắn gọn. Vui lòng
xem trong source code dự án)}

\hypertarget{sitemapspider-implementation-crawlerspidersspider.py}{%
\subsection{SitemapSpider Implementation
(crawler/spiders/spider.py)}}

\emph{(Mã nguồn chi tiết đã được lược bỏ để báo cáo ngắn gọn. Vui lòng
xem trong source code dự án)}

\hypertarget{kafka-pipeline-crawlerpipelines.py}{%
\subsection{Kafka Pipeline
(crawler/pipelines.py)}}

\emph{(Mã nguồn chi tiết đã được lược bỏ để báo cáo ngắn gọn. Vui lòng
xem trong source code dự án)}

\hypertarget{chux1ea1y-crawler}{%
\subsection{Chạy Crawler}}

\hypertarget{local-development-vux1edbi-docker-compose}{%
\subsubsection{Local Development (với Docker
Compose)}}

\emph{(Mã nguồn chi tiết đã được lược bỏ để báo cáo ngắn gọn. Vui lòng
xem trong source code dự án)}

\hypertarget{output-mong-ux111ux1ee3i}{%
\subsubsection{Output mong đợi}}

\begin{verbatim}
2024-01-15 10:30:45 [INFO] Spider opened: news_sitemap
2024-01-15 10:30:46 [INFO] Crawling VnExpress sitemap: https://vnexpress.net/sitemap_news.xml
2024-01-15 10:30:47 [INFO] Found 150 URLs in sitemap
2024-01-15 10:30:48 [INFO] Parsing article: https://vnexpress.net/kinh-doanh/...
2024-01-15 10:30:49 [INFO] Sent to news_raw[0]@123
...
2024-01-15 10:45:30 [INFO] Closing spider (finished)
2024-01-15 10:45:30 [INFO] Spider closed: finished
\end{verbatim}

\hypertarget{xuxe1c-minh-kafka-messages}{%
\subsubsection{Xác minh Kafka
Messages}}

\begin{Shaded}
\begin{Highlighting}[]
\CommentTok{\# Consume messages}
\ExtensionTok{docker}\NormalTok{ exec newsrag{-}kafka kafka{-}console{-}consumer }\DataTypeTok{\textbackslash{}}
  \AttributeTok{{-}{-}bootstrap{-}server}\NormalTok{ localhost:9092 }\DataTypeTok{\textbackslash{}}
  \AttributeTok{{-}{-}topic}\NormalTok{ news\_raw }\DataTypeTok{\textbackslash{}}
  \AttributeTok{{-}{-}from{-}beginning} \DataTypeTok{\textbackslash{}}
  \AttributeTok{{-}{-}max{-}messages}\NormalTok{ 5}
\end{Highlighting}
\end{Shaded}

Sample message:

\emph{(Mã nguồn chi tiết đã được lược bỏ để báo cáo ngắn gọn. Vui lòng
xem trong source code dự án)}

\hypertarget{dockerfile-cho-fargate-multi-stage}{%
\subsection{Dockerfile cho Fargate
(Multi-stage)}}

\emph{(Mã nguồn chi tiết đã được lược bỏ để báo cáo ngắn gọn. Vui lòng
xem trong source code dự án)}

\hypertarget{build-test-cux1ee5c-bux1ed9}{%
\subsection{Build \& Test cục bộ}}

\begin{Shaded}
\begin{Highlighting}[]
\CommentTok{\# Build image}
\ExtensionTok{docker}\NormalTok{ build }\AttributeTok{{-}t}\NormalTok{ news{-}crawler:local .}

\CommentTok{\# Chạy container (cần Kafka, PostgreSQL chạy)}
\ExtensionTok{docker}\NormalTok{ run }\AttributeTok{{-}{-}rm} \DataTypeTok{\textbackslash{}}
  \AttributeTok{{-}{-}network}\NormalTok{ host }\DataTypeTok{\textbackslash{}}
  \AttributeTok{{-}e}\NormalTok{ KAFKA\_BOOTSTRAP\_SERVERS=localhost:9092 }\DataTypeTok{\textbackslash{}}
  \AttributeTok{{-}e}\NormalTok{ DB\_HOST=localhost }\DataTypeTok{\textbackslash{}}
\NormalTok{  news{-}crawler:local python main.py }\AttributeTok{{-}{-}mode}\NormalTok{ crawl}
\end{Highlighting}
\end{Shaded}

\hypertarget{thuxeam-trang-tin-tux1ee9c-mux1edbi}{%
\subsection{Thêm trang tin tức
mới}}

\begin{enumerate}
\def\labelenumi{\arabic{enumi}.}
\tightlist
\item
  \textbf{Phân tích trang} --- Tìm \texttt{sitemap\_news.xml} (thường ở
  \texttt{/sitemap\_news.xml} hoặc \texttt{/sitemap.xml})
\item
  \textbf{Kiểm tra trang bài viết} --- Dùng browser DevTools tìm XPath
  selectors cho title, content, author, date, category
\item
  \textbf{Thêm vào config\_site.json} --- Copy entry hiện có và sửa
  selectors
\item
  \textbf{Test cục bộ} --- Chạy spider và xác minh trích xuất
\item
  \textbf{Deploy} --- Rebuild Docker image, push lên ECR
\end{enumerate}

\hypertarget{vux1ea5n-ux111ux1ec1-thux1b0ux1eddng-gux1eb7p-khux1eafc-phux1ee5c}{%
\subsection{Vấn đề thường gặp \& Khắc
phục}}

\begingroup
\small
\setlength{\tabcolsep}{3pt}
\renewcommand{\arraystretch}{1.15}
\begin{longtable}[]{@{}
  >{\raggedright\arraybackslash}p{(\columnwidth - 2\tabcolsep) * \real{0.4211}}
  >{\raggedright\arraybackslash}p{(\columnwidth - 2\tabcolsep) * \real{0.5789}}@{}}
\toprule\noalign{}
\begin{minipage}[b]{\linewidth}\raggedright
Vấn đề
\end{minipage} & \begin{minipage}[b]{\linewidth}\raggedright
Giải pháp
\end{minipage} \\
\midrule\noalign{}
\endhead
\bottomrule\noalign{}
\endlastfoot
\texttt{403\ Forbidden} & Thêm headers phù hợp, tuân thủ
\texttt{ROBOTSTXT\_OBEY}, giảm \texttt{CONCURRENT\_REQUESTS} \\
\texttt{XPath\ returns\ empty} & Kiểm tra site có dùng JavaScript
rendering (cần Splash/Playwright) \\
\texttt{Kafka\ connection\ failed} & Đảm bảo Kafka healthy, check
\texttt{bootstrap.servers} \\
\texttt{Memory\ error} & Giảm \texttt{CLOSESPIDER\_ITEMCOUNT}, tăng
container memory \\
\texttt{Timeout} & Tăng \texttt{DOWNLOAD\_TIMEOUT},
\texttt{CLOSESPIDER\_TIMEOUT} \\
\end{longtable}
\endgroup

\hypertarget{bux1b0ux1edbc-tiux1ebfp-theo}{%
\subsection{Bước tiếp theo}}

Sau khi crawler hoạt động cục bộ: 1. \textbf{Build \& Push lên ECR} ---
Xem \href{5.9-AWS-Deploy/}{AWS Deployment Prep} 2. \textbf{Đăng ký Task
Definition} --- Xem \href{5.10-Fargate-Crawler/}{Fargate Crawler} 3.
\textbf{Cấu hình SQS} --- Thay Kafka bằng SQS cho AWS deployment (xem
\href{5.11-Lambda-Consumer/}{Lambda Consumer})

\begin{center}\rule{0.5\linewidth}{0.5pt}\end{center}

\textbf{Tiếp theo:} \href{5.6-Ingestion/}{Nhập dữ liệu (Kafka Consumer
\(\rightarrow\) PostgreSQL)}
